\section{Benchmark 解读与研究方向}

\subsection{Arc AGI / O3 曲线解读}
Arc AGI 作为 reasoning benchmark 的核心指标，最早能让 1,000 美元预算达到 87.5\% accuracy，而 20 美元以下只能维持在 25\%。Chen 指出这条曲线并非只是“给钱就能做”，而是要同步排布 sampler、search 和 verifier，才能让每一笔预算都转化为有效 accuracy。OpenAI O3、AlphaProof/AlphaGeometry 等 demo 的实践，牢牢依赖 inference-time compute，而不是训练域。

\begin{table}[H]
\centering
\caption{主要 Benchmark 的推理时预算/性能对比}
\begin{tabular}{p{0.30\textwidth}p{0.28\textwidth}p{0.32\textwidth}}
\toprule
Benchmark & Token / Cost budget & Observed behavior \\
\midrule
Arc AGI & \$20~\$1,000 per query & 25\% → 87.5\% accuracy; the curve saturates after \$500 without better search/verifier \\
OpenAI O3 & \$20~\$200 & Consistency + trace feedback gives human-level accuracy once budget crosses \$100 \\
AlphaProof / AlphaGeometry & 50~200 token steps & RL-based reward tweaks inference-time policy to maintain correctness under limited latency \\
\bottomrule
\end{tabular}
\end{table}

\subsection{Benchmark 仪表与预算分层}
Chen 建议为每个 benchmark 都建立一个 ``budget bin'' 仪表板：把 token/cost 预算分为 Conservative、Balanced、Aggressive 三档，分别对应 `20~80 token`、`80~300 token`、`>300 token` 的情况。每个 bin 都要记录 accuracy、stepwise verifier margin、token per vote，以便在运行时快速定位“当前状态是哪个 bin”。这种分层可帮助团队在对外发布 benchmark 结果时直接说“在 Balanced bin 下，我们的 accuracy 87.5\% 但 latency 3.4s”。

\begin{table}[H]
\centering
\caption{Benchmark 仪表示例}
\begin{tabular}{p{0.26\textwidth}p{0.30\textwidth}p{0.30\textwidth}}
\toprule
Bin & 典型配置 & 监控指标 \\
\midrule
Conservative & $\leq 80$ token, $\leq 2$ iterations & Accuracy, sampling diversity, latency \\
Balanced & 80~300 token, 2~3 iterations & Verifier margin, branching ratio, token per candidate \\
Aggressive & > 300 token, 3+ iterations & Score per token, trace iterations, SLO breach rate \\
\bottomrule
\end{tabular}
\end{table}

\begin{importantbox}{Benchmark 仪表的落地建议}
1. 为每个 bin 配置一个 `alert level`，当 latency > target 或 verifier margin < threshold 时切换 bin；2. 把 bin 的 state push 到 dashboard，让运营团队可以快速阅读“今天的 benchmark 运行在哪一档”；3. 结合 bin 内的 score-per-token 监控，及时发现 token budget 失控的 early warning。
\end{importantbox}

\subsection{开放研究问题}
\begin{itemize}
\item 如何在低 latency 场景下用有限 compute 拓展 search depth 与 iteration depth？
\item 能否设计一个“budget-aware verifier”实时报告 token/candidate tradeoff，而不等到 inference 完成？
\item Trace feedback 在多模态/图像+文本 reasoning 中的通用性如何保障？
\item 在真实部署 pipeline 中如何把 sampling diversity、stepwise scaling metric、trace drift 等指标串成可执行的 SLO？
\end{itemize}

\begin{itemize}
\item 推理时 pipeline 能否用 learning-to-rank 的方式自适应调整 branching ratio 与 iteration depth，而不是人工设定阈值？
\item 对于 adversarial prompt，verifier margin 如何与 prompt-editing loop 结合，以免错误推理被 feedback 强化？
\item 多模态 reasoning 下，token 预算与 vision token length 的对应关系还有哪些不确定性？
\end{itemize}

这些问题在 01:19:00 之后的 Q\&A 被反复提及，Chen 建议通过 “metric-first hypothesis” 方式探索：先提出核心 metric（如 `score per token`）再迭代调整 pipeline，而非凭直觉扩大 token budget。

\begin{knowledgebox}{研究方向联想}
开放问题可以从 budgeting instrumentation、verifier robustness 以及 trace 科学三个角度切入。Chen 建议把 inference-time logs（sampling count/delay/score）写入 dashboard，为未来自动调整策略提供基础数据。
\end{knowledgebox}

\subsection{本章小结}
通过对 Arc AGI、O3 以及 DeepMind 系列的 budget/accuracy 曲线解读，我们看到 inference-time compute 既是工程参数也是研究方向的触发器。未来需要在 low-latency 下维持 high accuracy，并把 trace/verifier/gating 指标织入部署 pipeline。

